\section{Tradition and Being}

\begin{quotex}
We think that if a Western tradition could be rebuilt it would be bound to take on a religious form in the strictest sense of this word, and that this form could only be Christian; for on the one hand the other possible forms have been too long foreign to the Western mentality, and on the other it is only in Christianity---and we can say still more definitely in Catholicism---that such remnants of a traditional spirit as still exist in the West are to be found. Every `traditionalist' venture that ignores this fact is without foundation and therefore inevitably doomed to failure; it is self-evident that one can build only upon something that has a real existence, and that where there is lack of continuity, any reconstruction must be artificial and cannot endure.
\flright{\textsc{Rene Guenon}, \emph{The Crisis of the Modern World}}
\end{quotex}


As we get close to the end, it is good to emphasize Guenon's point yet again. The whole purpose of Gornahoor is to go back down to the foundations of the Western Tradition to see how it can possibly be rebuilt. This actually involves two interrelated task:

\begin{enumerate}


\item Going beyond the outward form of the Tradition to recover its inner meaning

\item Developing a real knowledge of Eastern doctrines
\end{enumerate}


The second task serves as a corrective to the first. In other words, Eastern doctrines may serve as the key to extract the inner meanings from the first task. Moreover, ``foundations'' include the ancient Greek, Roman, Nordic, and Celtic elements that had been assimilated into Christianity. In particular, the Middle Ages, being closest to us in space and time, even if different in terms of intellectuality, are vital to this process of rebuilding. Although there is considerable resistance to this project, even from those who consider themselves ``traditional'', Guenon is hardly sanguine about alternatives. He writes:

\begin{quotex}
The least fantastic venture, in fact the only one that does not come up against immediate impossibilities, would therefore be an attempt to restore something comparable to what existed in the Middle Ages, with the differences demanded by modifications in the circumstances; and for all that has been completely lost in the West, it would be necessary to draw upon the traditions that have been preserved in their entirety, as we stated above, and, having done so, to undertake the task of adaptation, which could be the work only of a powerfully established intellectual elite. All this we have said before, but it is useful to insist on it again because too many inconsistent fantasies are given free rein at present.
\end{quotex}


Some of those fantasies, and objections, can be described as follows.

\subsubsection{Confusion with Church Organization}

Catholicism, as the Western Tradition, stands on its own. Besides Holy Scriptures, the place to find it is in the writings of the Fathers, Doctors, and Saints of the Church, inter alia. One of the roles of the Bishops is to preach and hand down that tradition. Some may perform that role better than others. However, they do not ``own'' the Tradition. All too many people presume that Catholicism is just about the church organization. If a bishop is mistaken, or acts badly, that means nothing for understanding the Tradition.

\subsubsection{Presentism}


Catholicism has very deep roots and cannot be limited to the acts and opinions of Church leaders today. That is not what we, nor Guenon, are referring to. On the contrary, we should not be surprised since Guenon explicitly denies that contemporary church leadership has an in depth understanding of Tradition.

\subsubsection{Confusion about Orthodoxy}

From the point of view of the esoteric Tradition, Catholicism refers also to Orthodoxy. That is a single Tradition, even if there are now competing church structures on the exoteric level. A heresy is a break from tradition, but a schism is not.

\subsubsection{Syncretism and Indifferentism}

Guenon is not advocating syncretism, i.e., the arbitrary combining of elements from disparate traditions. Nor should one be indifferent to his own Tradition. Nevertheless, at the true intellectual level there can be fruitful interaction. This is nothing new. Gornahoor has amply documented how pagan Greek ideas influenced the Fathers. The Schoolmen were quite willing to engage in dialogue with Jewish, Muslim, and Pagan thinkers.

\subsubsection{Paganism}


There is no religion called ``paganism'', so a return to it is impossible. What it means today is ``non-Christian'' or more usually, ``anti-Christian''. Such movements are useless, fantastic ventures, with no intellectual depth. If you fantasize about creating a post-Christian Tradition, then you would still need to pick and choose those elements from the Christian Tradition to assimilate.

\subsubsection{Turning East}


Engaging with Eastern doctrines is not the same thing as ``converting'' to the outward forms of an Eastern religion. In almost all cases, this means adhering to the religion of your forefathers. It is a bourgeois conceit to believe that you can change your religion, something that would have made as little sense to an ancient as changing one's sex or race would have.

\paragraph{False Traditionalism}


\begin{quotex} ``traditionalists'' refer to people who only have a sort of tendency or aspiration towards tradition \emph{without really knowing anything at all about it}; this is the measure of the distance dividing the ``traditionalist'' spirit from the truly traditional spirit, \emph{for the latter implies a real knowledge}... 
\end{quotex}


Let Guenon's words really sink in. There are many who call themselves ``Traditionalists'', as if it were no more than one philosophical, or even worse ``political'', movement among many. Someone may say he ``converted'' to Traditionalism, or from Traditionalism to some other religion. That is impossible, like a chef saying he converted from reading cookbooks to actually cooking food. Such men know nothing at all, so stay away from them.

Tradition is the normal state of human culture, in almost all times and places. It is the ``civilization of space'', as Evola put it in \textit{Civilizations of Space and Civilizations of Time}\footnote{\url{https://gornahoor.net/?p=12696}}. There is only the question of recognizing it or not.

Hence, the normal situation is to follow a Tradition; that is what makes a ``traditionalist''. It is a state of being, not a system of thought. At some point he may discover Guenon, which helps him understand his own Tradition in much more depth.

\paragraph{Forgetting Being}


Heidegger made one good point about the forgetfulness of Being. That means we have exchanged the idea of something for the thing itself. For example, I've known many who are enamored with the idea of being a Buddhist, but few, if any, who strive to purify the mind to perfectly reflect the Buddha nature. As Dogen pointed out, just calling yourself a Buddhist does not lead to enlightenment without proper practice.

This forgetfulness is the result of embracing nominalism rather than realism. For the nominalist, what matters is the ``name'' of things. A Buddhist is just someone who claims to be. However, the one thing needful is a real change in Being. That requires an inner transformation from a state of less being to more being. Specifically, that means the actualization of one's possibilities. As you are now, your Real Self and True Will are only virtual and need to be realized.

\paragraph{Addiction to Thinking}


Like sex, food, or drugs, we can become addicted to thinking. It makes us unfree, since we live a factitious life in our mind that may have little contact with reality. The lower intellect is mechanical. It responds to an idea with either a ``like'' or a ``dislike''; for proof, just look at how popular facebook has become just on that notion.

It also works on Yes and No. If Tom says Yes to an idea, then Jerry feels compelled to say No and start an argument.

Thought is corrosive of systems whenever it sets itself up as the supreme judge. Yet, there is a knowing that transcends mere rational thinking. Is it better to be courageous or to know how to define it? Is it better to be just or to argue about it?

Purification of the mind is necessary to think well. It will never suffice to explain an idea clearly, logically, and with compelling evidence as long as ``the reason is perverted by passion, or evil habit, or an evil disposition of nature''.

Religion teaches us how to transcend our passions, evil habits, and bad dispositions. It is much more than a system of thought. One's nominal affiliation proves nothing. You can argue, if you like, about religious or metaphysical ideas. I, on the other hand, am more interested in whether you follow the (traditional) religion of your fathers, have a male priesthood, and pray for the dead. That is Being.


\flrightit{Posted on 2016-04-28 by Cologero}

\begin{center}* * *\end{center}

\begin{footnotesize}
\begin{sffamily}
\texttt{Johannes on 2016-04-29 at 01:26 said:}

Conversion to catholicism is probably the worst advice anyone can give.

\hfill

\texttt{Cologero on 2016-04-29 at 03:11 said:}

Johannes, is that the conclusion you reached? This is what I invited the reader to do:

\begin{itemize}


\item Go beyond the outward form of the Tradition to recover its inner meaning

\item Develop a real knowledge of Eastern doctrines
\end{itemize}


Did no one read that?

\hfill

\texttt{B on 2016-04-29 at 08:41 said:}

Many of the points here is what myself and close others have come to as the conclusion... to choose a living established European tradition, and stick with it through and through. The nice thing about the Catholic Faith is that it proves itself True quickly the more you learn.

\hfill

\texttt{Adam Wallace on 2016-04-29 at 08:48 said:}

\emph{``Engaging with Eastern doctrines is not the same thing as ``converting'' to the outward forms of an Eastern religion. In almost all cases, this means adhering to the religion of your forefathers. It is a bourgeois conceit to believe that you can change your religion, something that would have made as little sense to an ancient as changing one's sex or race would have.''}

What would be the wisest course of action for a someone who is born into a totally secular environment, without any father figures of any religiousity? Following one's nature, so to speak? The adoption of any residual tradition in one's vicinity, no matter how watered-down? Or is such a person essentially doomed to his circumstance?

\hfill

\texttt{Lu Dongbin on 2016-04-29 at 09:39 said:}

Many of the Eastern doctrines stress the need to meet an enlightened master in order to actually succeed in attaining liberation. This is why Chan Buddhist monks and Daoist adepts would travel from mountain to mountain and across the realm in search of a true master.

As far as I can see it, modern Catholicism and Christendom in general lacks enlightened masters and increasingly is entirely lacking in vitality. For one sincerely searching for transformation and gnosis, what can modern Catholicism provide to assist in this endeavor? One may be able to make progress trying to implement the wisdom and discipline suggested by the Desert Fathers or the medievals, but without an initiatory encounter with a true master can one truly hope to attain transcendence? I am not saying one will necessarily find this in the East either, but it seems much more likely.

Furthermore, you say no one can change their religion. What of those born without a religion, who grew up in secular households and were without a tradition from birth?Those may also be a product of the modern, multicultural West and have grown up as much around Eastern traditions as anything else. In such a case is choosing one's ancestral religion, which may be divorced from such a person by three generations or more, the only logical choice?

Finally, for all of Guenon's solid reasoning here, why did he become a Sufi Muslim? Could it be he saw that modern Christianity no longer offers the possibility to achieve what he sought to achieve?

Don't get me wrong here, I am supportive of Christianity and would love to see a true restoration as described. That said, excluding the hippy New Age types, I can see why sincere Western seekers may turn East in the present circumstances.

\hfill

\texttt{Scardanelli on 2016-04-29 at 10:45 said:}

It seems that for every truly clarifying post, 10 comments arise that seek to muddy the waters.

Lu Dongbin:

When did Cologero advise you to practice ``modern Catholicism''? Did he not, in fact, advise the complete opposite? Catholicism in fact grants you access to not just to A master but THE spiritual master in Christ. Follow the instructions of the post and encounter the tradition in the writing of the saints, desert fathers, and doctors in truth and fact rather than merely theoretically. Pray, fast, meditate, and learn what is meant by ``concentration without effort'' through experience of that state, rather than merely thinking or writing about it. The path begins in ``the heart'' not the brain. Then ask questions.

As to your second point, I would go so far as to say that anyone born in the west is born as a Christian. The outward conditions are irrelevant. The west is The Catholic tradition and the Catholic tradiation is the West. There is no one without the other as Cologero has written many times. You were born outside of an exoteric form... Fine. That does not exclude you from experiencing the inner truth of the tradition.

\hfill

\texttt{JB on 2016-04-29 at 11:01 said:}

What Rene Guenon wrote publicly in his books on Christianity and the future of the West clashes singularly with what he wrote privately in his correspondence, Cologero, and you know that, so to pretend otherwise and mislead your readers is bad faith on your part.

Your blog does not address the christianisation of European peoples, with respect to the notion of conversion, and I wonder why. Could it be that it contradicts what you write in the next sentence?

As for conversion to other traditions, you ought to know by now that certain other traditions do not recognize ``conversion'' in the way it is generally understood in the christian framework. If you are still not familiar with this take on things, more fool you.

\hfill

\texttt{Tom on 2016-04-29 at 11:52 said:}

Let me be another of the Jerry commenters saying no to Tom's Yes. I admit that the following rebuttals are cliched but I feel they should be repeated since they are pertinent.

``If you fantasize about creating a post-Christian Tradition, then you would still need to pick and choose those elements from the Christian Tradition to assimilate''

What is wrong with that though? You say yourself that the Scoolmen were willing to engage with pagan thinkers and the pagan elements of Christian Tradition are well known. Islam and Christianity each started, not with an unbroken tradition, but with a break from a tradition, the former actively attacking the sacred relics of its predecessor in the Arabian world.

Any such criticism of those who seek to revive a polytheistic spiritual practice without an unbroken line of Tradition could equally be applied to the founders of any religion.

The intention, in every case, is to get to the truth, which has always been the same

\hfill

\texttt{Cologero on 2016-04-29 at 13:03 said:}

@JB, I don't know what Guenon wrote privately, but it is his public record that ultimately counts. He could have retracted had he wanted to.

If you are able to get past your inner rage, you may have noticed that I did address the question of a different future for the West; you could have commented on that if only you had the requisite intelligence. Of course, if you insist on following Guenon, a return to some pre-Christian tradition is out of the question. He also claimed that no new tradition is possible. So what do you suggest as a replacement, other than secularism? As you know, Guenon speculated that it would be Islam.

We certainly have addressed the ``christianisation'' of the European people. It was a national decision. For example, you can read about the conversion of Iceland in \textit{Spiritual Invigoration}\footnote{\url{http://www.gornahoor.net/?p=8015}}.

Your comment about ``conversion'' is pointless. It means to apostasize from one religion and adopt another. Are you not familiar with that?

In the future, I would ask all commenters to make sure their Xanax kicks in before clicking the ``Reply'' button.

\hfill

\texttt{Lu Dongbin on 2016-04-29 at 13:25 said:}

@Scardanelli

As stated, he also said to integrate Eastern wisdom. Much of the Eastern wisdom I've read emphasizes how vital it is to encounter a true teacher who can show you the subtleties of the way. If a tradition loses this element it seems to me such a tradition has grown old, weary, and ineffectual. We may be able to ignore the lack of such a master in modern Christendom and practice what the ancients left us, but while such will likely lead to purification and progress on the path, certainly superior results to doing nothing or being a secularist like everyone else, will it lead the most serious of seekers to transcendence, to gnosis?

If the visible, living form of a tradition is utterly without vitality and indeed for the most part consciously promoting the most poisonous elements of modernity and we need to invoke an invisible, non-human teacher to guide us, aren't we in a similar position to the neopagan? The medieval tradition is gone in the same way ancient Neoplatonism is gone--we have the texts and practices of both, so where does the difference lie? In that the empty shell of the organized form lives on for the former?

If Catholicism is the West and the West is Catholicism, a notion I generally agree with, then it seems to me that the West is dead. Indeed all signs around us point to such being the case. Hopefully out of the ashes something new will arise, I just hope it won't be Islam for a variety of reasons.

\hfill

\texttt{Logres on 2016-04-29 at 18:39 said:}

Adam W, Some time ago, when wandering in the morass of Western thought, wondering what to make of the remains of the day, knowing at best just pieces of Tradition (in the form of scholarly trails and remarks), I started praying to Saint Alban to show me in what direction I should move. St. Alban lost his head literally, and became one of the patron saints of England, from where most of my ancestors came. I have to say that that prayer was answered beyond what I could have imagined. Even if you have no physical ancestors to teach you, there are spiritual ones.

\hfill

\texttt{Harun on 2016-04-29 at 18:46 said:}

Christianity lacks spiritual teachers, this is true but not completely.

The Church is now dominated by modernists and is dying together with the West who was once christian, but on the inside of the Church lie something that can never be destroyed.

The Apostle John will wait until the end of times, preserving true gnosis, there are few in the Church who will always preserve Tradition ``and the gates of hell shall not prevail against it''.

This is the final part of Kali Yuga, Christianity is not a traditional form among others that will die and be replaced like the Roman religion 2000 years ago. Other Traditions are following the degeneration fo the Church, it merely started in the West and so the Church was corrupted first.

Guénon wrote that just a few people are necessary to preserve a Tradition, even if it's unknown to the world.

So it will be in the last days, the last keepers of all Traditions will be united when the Kali Yuga ends.

Despite the lack of ``gurus'' with expericence of theosis in modern christianity, its spiritual influence is not dead.

Theosis inside christianity is still virtually possible and it happen even today, even without a ``guru''.

Isn't Padre Pio a fully realized ascetic with extraodinary siddhi?

De Giorgio visited him and think he was, Guénon was very interested.

I also know Padre Pio visited in a dream an italian hermeticist and foresee his death saying ``you are a creature of heaven and you'll return soon to heaven''.

What about Mount Athos? St. Silouan certainly reached the highest form of hesychasm and he taught many young monks, and this merely 80 years ago. Some of his disciples died in the 90's and taught many other monk who are still alive today.

Not to forget the secret ``legend'' about the twelve invisible hermits of Mt.Athos: it's a symbol of a very powerful spiritual influence who still exist today.

Do not mistake rare for impossible.

\hfill

\texttt{JB on 2016-04-29 at 23:28 said:}

Here are two quotes from Guenon's correspondence which perfectly illustrate what I am talking about:

~ «~Pour ce que j'ai dit dans Orient et Occident au sujet du rôle possible de l'Eglise catholique (comme représentant une forme traditionnelle occidentale pouvant servir de base à certaines réalisations, ainsi que cela a d'ailleurs eu lieu au moyen âge), je dois dire que je ne me suis jamais fait d'illusions sur ce qu'il pouvait en résulter en fait des circonstances actuelles~; mais il ne fallait pas qu'on puisse me reprocher d'avoir pu négliger certaines possibilités, au moins théoriques, ou de ne pas en tenir compte.~» (Lettre à Roger Maridort, 1933).

«~Vous devez bien penser que je ne suis pas si naïf que cela~; mais pour des raisons qu'il ne m'est malheureusement pas possible d'expliquer par lettre, il était nécessaire de dire ce que j'ai dit et d'envisager cette possibilité, ne fût-ce que pour établir une situation nette, et a eu pleinement le résultat (négatif) que j'attendais.~» (Lettre à Julius Evola, 1933).

You can read an article on the subject (where those quotes were found) here:

\url{http://www.la-question.net/archive/2009/09/04/rene-guenon-un-esoteriste-antichretien.html}


On conversion, if we are to understood you correctly, the ancient Europeans somehow suddenly all developed ``bourgeois conceits'' and converted to an exotic foreign religion, instead of following the faith of their fathers. How do you justify this?

Your article on the conversion of Iceland is hardly representative of Europe as a whole. Contrasted with the christianisation of Rome, followed by the massacres of Charlemagne, there is a world of difference, both in rationale and method.

To conclude, to rely exclusively on Guenon, a man who was more influenced by nineteenth century occultism than by a genuine knowledge of traditional doctrines, contrary to what he continually maintained, is not only misguided but foolish.

\hfill

\texttt{Nilakantha on 2016-04-29 at 23:30 said:}

There are contingent reasons that necessitate changing one's traditional affiliation. As Guénon himself said: I t too often happens that someone is born into a milieu not in harmony with his own nature, and because not really suitable for him, does not allow his possibilities, especially of the intellectual and spiritual order, to develop in a normal manner. Certainly it is regrettable in more than one respect that things are this way, but these are the inevitable drawbacks of the present phase of the KALI-YUGA.

\hfill

\texttt{Cologero on 2016-04-29 at 23:53 said:}

Close, JB, so what's the point? You both consider Guenon as an authority and then not as an authority. But to comment on the content of that correspondence ... He is referring to the ``Catholic Church'', not to Catholicism as a Tradition. We anticipated that objection if only you could read English as well as your French. Actually, we bemoaned that fact in several texts. Specifically, we pointed out how Guenon was not able to make much headway with the Church, and was particularly opposed by Jacques Maritain, unfortunately. (You can do a search for them.) That was the basis of his complaint to Evola.

Nevertheless, Guenon honored Pope Pius XI's condemnation of Action francaise on the grounds that he was the spiritual authority.

So the conclusion is not that Guenon was anti-Christian or anti-Catholic, but rather that he did not believe that the organized Church was still capable of reviving the very Tradition of which she was the visible embodiment. That is simply a prudential judgment, applicable to a specific time and place. Guenon acted on that judgment and remained in Egypt where his adopted religion was the right fit.

On several occasions, we have offered another judgment, based on the publications of some remarkable books by Valentin Tomberg an Boris Mouravieff. They could not have appeared if the Tradition was totally lost. They reveal the Church of John is still active. Since we've dealt with that in some depth, I won't repeat it here. So we have carried on their work ... you can take it or leave it, as you see fit.

JB, you could have saved yourself the risk of carpal tunnel syndrome by stopping after the two quotes. The last 4 paragraphs are incoherent, so I can't comment on them. But if you are seeking my approval to change either your religion or your sex, my advice is only to follow your informed conscience.

\hfill

\texttt{clay masterson on 2016-04-30 at 11:09 said:}

Nostalgia for the past is a surefire road to failure, we are evolving beings, not stagnant. New forms and new traditions are needed, not a retreat to the past.

\hfill

\texttt{Cologero on 2016-04-30 at 12:14 said:}

Mr Masterson, you have managed to combine so many errors, misconceptions, and misreadings into two brief sentences, so I can only conclude it is not the result of chance. In my moments of weakness, however, I have to wonder what it is that you have evolved into and what new tradition you follow!

\hfill

\texttt{Harun on 2016-04-30 at 15:34 said:}

``The spiritual destinies of such men, the paths they take, are only determined by the orientation of their own nature and the `calls' from purely spiritual domains that are `heard'. Their decision is rooted in such inner factors and may only be understood by insight into their inner nature and how it relates to the divine order''

I can relate to this.

I met a qualified representant of Buddhism but he's teaching sunyavada, and inside me everything told me it was not my path.

The problem is if it's wise to reject an initiatic path while we are still unable to follow an alternative path?

If we look at the life of Guénon I think it's the same decision he made when he waited many years in Europe and later became a Sufi instead of moving to India.

\texttt{Marcel K. on 2016-05-01 at 12:36 said:}

A. K. Coomaraswamy's son has written a detailed account of the change that took place within the curia when a dedicated group of Cardinals effected a turn towards liberal-humanist inspired ideas. He mentions that at a time a Marxist intellectual in rather well-known Italian feuilleton literally commented that ``the devil had entered the Council''. The conservative South Italian members of the curia were wholly unprepared for the onslaught of the new ideological current that emphasized the mundane `human warmth of Christianity' to the detriment of the traditional teachings. The `battle' was lost long before it began, because unlike their opponents the conservative members of the curia were at the grave disadvantages of having no particular goals other than the adherence to the traditional framework and were therefore quickly branded mere `obstacles to progress'.

If it was so easy for the democratic-profane forces to effect a virtual revolution within the Church, the pessimism of many so called traditionalist regarding Catholicism would seem quite arrogant. Rather than abandoning it, one should attempt to change the Church from within -- the Church is the domain of traditionalists after all, so with some shrewdness it should not be too difficult to expel the intrusive profane forces which quite evidently do not belong in it.

\hfill

\texttt{Nilakantha on 2016-05-01 at 13:37 said:}

The mythic Jesus worshipped by Christians as the plenary manifestation of God is not the Jewish messiah. The Jewish messiah has a political function intimately tied up with the ancient Kingdom of Israel. Jesus, on the other hand, is the bearer of a universal revelation formed for and by the Roman civilization centred in the Mediterranean basin, which is why Jesus is God for Europe.

Krishna, along with the other avataras of Vishnu, is a bearer of God's revelation for greater India. Whether or not you can ``become'' a Hindu is a question that must be honestly answered before you can authentically become a worshipper of Krishna. Regardless, we should never speak badly of any form that God has used in any orthodox religion to call men back to the truth.

\hfill

\texttt{Nilakantha on 2016-05-01 at 20:03 said:}

We cannot judge the orthodoxy of any religion by the miracles it produces. We must, on the other hand, judge the source of miraculous phenomena by the orthodoxy of the religion that produces them. We judge a religion orthodox because of its metaphysical doctrine, ethics and spiritual practice. Within an orthodox religion, miracles provide ``extrinsic and probable proofs'' of the spiritual efficacy of the tradition. Within a heterodox religion, miracles are merely demonic deceptions.

\hfill

\texttt{Lu Dongbin on 2016-05-02 at 22:27 said:}

I suppose for me the problem is that I am a bit of a pseudo-Traditionalist, insofar as I agree with most of what Guenon and Evola and company have to say, but I also can't seem to bring myself to agree that because potentially the major traditions can lead us to truth that they are therefore of equal quality or value. It seems to me that certain systems, ways of looking at the world, and forms of praxis are superior to others. In this regard I tend to see the ancient Aryan traditions, Neoplatonism, and Eastern traditions like Daoism and Chan/Zen as clearer, more potent, and overall superior in their ability to produce sages than the Abrahamic traditions. Of the latter, I respect traditional Catholicism the most, but I also believe it is due to the absorption of the European racial/pagan spirit and Neoplatonism that medieval civilization was glorious. Outside of Sufism and perhaps some of the insights of the Kabbalah, I have little interest or respect for either Islam or Judaism. This way of viewing things may reflect immaturity on my part and put me outside of some of the considerations of the Traditionalists, but that's how I see it.

I also cannot bring myself to accept exoteric dogmas that I know to be false. The prime issue barring me from Christianity is that I remain convinced there is an unborn, uncreated, timeless element at the core of man. This is heresy to Christianity and historically those who vocalized this heresy (like Eckhart) were condemned. (One of my other favorite Christians, Eriugena, was also condemned.) Of course I also cannot accept the notion of eternal damnation, divine linear time, animal souls dissolving at death, and other dogmas of Christianity. Should I ignore what I know to be false simply to practice the tradition of my ancestors, or perhaps I should say my most recent ancestors? One could retort that actually going through the transformation offered by discipline and practice is most vital, but why should deceiving oneself be required in the process? Combine this with the wretched state of modern Christianity and I see no reason to forgo the chance offered by this precious human life for transcendence in consideration of upholding my civilization's traditional religious form.

So for me, the choice has been to turn East--if this is a bourgeois conceit so be it. That said, I am not one to dispose of my heritage and try to become an Easterner like the New Age/hippie crowd. I try to remain connected to my racial traditions and culture insofar as possible in the modern world and, despite my own inability to follow the Christian path, I continue to support traditionalism in a Catholic context and hope for a revival in the Church.

\hfill

\texttt{Nilakantha on 2016-05-02 at 23:22 said:}

There is an absolute necessity to practice an orthodox religion, even though there's no need to take the myth of the religion literally. As Michael Griffin explains in the introduction to his Olympiodrus: Life of Plato and On Plato First Alcibiades 1-9: Olympiodorus treats the language of Christian doctrine, when it is mentioned, as if it belongs to the exegetical `level' of myth or phantasia. Like any myth, however, it is not simply false: it refers to real facts within the human soul or psyche. We may speak of Hera, or a certain `power' of the Christian God, but in either case the philosopher will recognise that what is meant is a certain capacity of the soul. We may speak of daimons or angels, but in both cases we are genuinely referring to human conscience. The overt language used, whether pagan or Christian, lies at the level of phainomena or appearances. Thus Olympiodorus is able to represent `philosophy' as a more fundamental framework that can accommodate various mythological systems and worldviews and facilitate their agreement.

\hfill

\texttt{Nilakantha on 2016-05-03 at 00:51 said:}

Bryan,

Your inability to find an authentic spiritual director is a problem all serious practitioners experience at this time. The best advice to give is also the most basic: be patient. God is the ultimate guru, and, like it or not, you have to learn to be patient and wait on His will to make itself known. Pray, fast, go to confession and attend Mass as often as possible and have faith that God has not left us orphaned.

The worst thing to do would be to immerse yourself in some kind of spiritual Esperanto of your own divising. Trying to cobble together your own unique spiritual path by picking and choosing bits and pieces from various traditional doctrines will just end up proving the old Muslim adage: He who uses himself as his own spiritual director has Satan for his guide.

And lastly, read \textit{On Gurus and Spiritual Direction} by Rama Coomaraswamy, a Traditionalist Catholic:

\url{http://www.worldwisdom.com/public/viewpdf/default.aspx?article-title=On_Gurus_and_Spiritual_Direction_by_Rama_Coomaraswamy.pdf}


\hfill

\texttt{/me/ on 2016-05-03 at 01:04 said:}

``There is an absolute necessity to practice an orthodox religion''

I understand that this is Tradition 101. However, if there probably is no such thing as a stupid question in mundane matters, certainly here with such a vital and essential matter I feel justified in asking, ``why''?

\hfill

\texttt{Johannes on 2016-05-03 at 02:42 said:}

Cologero, I personally considered joining the Catholic church at the suggestion of this blog. I talked to several different priests,went to traditional mass a few times a week, read the mystics, medieval theology and prayed 5-15 decades of the rosary daily. However in the end I just couldn't do it; my heart said ``no, this isn't true''. I wasn't raised in that tradition and just didn't believe in it . To me, its obvious that Catholicism is a false religion, and I was just forcing myself into it, hoping to find some `tradition' for myself. Look at the silly ramblings of those female mystics, or how far Catholic theology has diverged since the schism between east and west, let alone since VC2. There too many absurd dogmas to believe in, like the co-redemptrix and papal infallibility, which were clearly just made up by the priests to suit their needs as they went along, not `revealed' by some transcendent `Tradition'. What thinking person would willingly submit to believe in all that? Just because some old priest said so? Believing in things because some authority said so isn't possible anymore. That's a premodern mindset and its impossible to just `go back there', no matter how many books we read.

It seems to me that your kind of reasoning towards being a Catholic is only possible for someone who was born into that tradition, and then left it later on, only to realise it was due to youthful mistakes/arrogance/pride or whatever. But for some one who has no experience of Catholicism, it is too exotic, with its hundreds of Saint days and strange rituals that don't have much explanation other than `tradition'. The Catholic church has clearly been subverted in some capacity since Vatican 2. Catholicism is whatever the priests say it is; if they want to protestantise it and embrace modernity, well that's just Catholicism now. It doesn't make sense to be a traditional catholic believing all that nonsense, let alone now in 2016 where the priests have jettisoned most of that stuff. Your `hidden esoteric dimension' of Catholicism only exists in your mind and books, Cologero.

Having experienced it, Catholicism has no resonance with me, just like with most people of my age group. The imagined glorious medieval Catholicism doesn't exist, probably it never did, just in your minds, you just don't see this because you have fond memories from your childhood about it. But now it's just some musty old mumbling and mummery that the average westerner has no connection to. I understand that you are at least 70 years old Cologero, so perhaps that is a factor in your case. You are old and were raised in a time when perhaps that religion had some cultural impact, at least on your life. But for young people today, it cannot be that way.

No doubt you will accuse me of `raging', but rage against what? I was never a Christian, I have nothing to rage against, no issue with Christians. If they want to believe in that well good for them. However I am annoyed at you personally for giving me stupid advice and wasting my time.

\hfill

\texttt{Matt on 2016-05-03 at 12:12 said:}

Johannes,

Co-redemptrix is not a dogma. It's a theological concept that some within the church accept and others do not. Maybe if more attention was paid to studying that religious tradition you had considered joining, you wouldn't have been left with such dissapointment.

\hfill

\texttt{Cassiodorus on 2016-05-03 at 16:18 said:}

Lu Dongbin,

I'm curious- what are the exoteric beliefs of Christianity that you cannot believe?

\hfill

\texttt{Cassiodorus on 2016-05-03 at 16:32 said:}

Lu Dongbin, in some sense, I agree with you. To my lights, Christianity and unqualified nondualism are not compatible. But, why should we be so quick to reject panentheism in either a weak or strong form?

\hfill

\texttt{Cologero on 2016-05-03 at 22:27 said:}

I have disappointed many people in my life, Johannes, mostly women however. Perhaps you all should create a facebook group where your juvenile whining will find a sympathetic ear.

I have striven to outline a jnana path that certainly has existed in the Tradition. That includes, at a minimum, a correct metaphysical worldview, a method of self-knowledge, and a path to self-mastery over one's thoughts. Have you done any of that? Have you looked into the many texts and practices that have been mentioned in these pages? There are three mailing lists in which such issues could be discussed. In addition, we have weekly training in self-knowledge exercises ... these are real, not texts in a book. You've never participated. And not just you, but all the others who are whining about not finding a teacher. A man takes responsibility for his own education, not to mention the consequences of the choices he has made.

I don't recall recommending all those things to you. For example, what women mystics are you referring to and where did I suggest them to you. I fully understand that a sentimental bhakti path if not very suitable to young men today. You make the mistake I warned you about: viz., confusing structures with the Tradition itself. So which metaphysical teachings do you reject? The unity of being? The cosmic order? The hierarchical organization of the universe?

Faith is a matter of Will. Had you actually paid attention to my advice, you would have learned inner silence, you would have gained authentic gnosis, you would have become a man of daring. Perhaps you are simply not up to the task.

\hfill

\texttt{Nilakantha on 2016-05-03 at 23:21 said:}

/me/,

The path from our everyday existence to the unconditioned requires three elements. First, we need a metaphysical doctrine to act as a map. But, since map is not territory, we also need concrete practices with which to make the journey. Our second element is an ethics to pry us away from our fascination with the multiplicity of daily life. As Porphyry says:''If we are to speak frankly, concealing nothing, there is no other way to achieve our end than by being riveted (so to speak) to the god, and unfastening the rivet of the body and the pleasurable emotions of the soul which come from the body; security comes to us by actions, not just by listening to lectures.It is not possible to be familiar with a god -- not even with one of the particular gods, let alone the god who singly is above all and higher than incorporeal nature -- by following just any lifestyle, especially flesh-eating; one can hardly, even with all kinds of purifications of soul and body, become worthy of awareness of the god, that is if one has a fine nature and lives a pure and holy life'' --- On Abstinence from Killing Animals. And he says in another place: ``It is impossible for a man who loves God also to love pleasure and the body, for he who loves these must needs be a lover of riches. And he who loves riches must be unrighteous. And the unrighteous man is impious towards God and his fathers, and transgresses against all men'' --- Porphyry's Letter to His Wife Marcella. The third element is a spiritual method by which the mind, having been freed to some degree from concern with the body concentrates itself on a symbol of the Absolute. This often involves the recitation of a divine name or mantra with the aid of a ma?la? or rosary. If they are to be effective, the three elements of orthodox religion must have been revealed by God, since He knows the path and the way to traverse it --- thus our need for an orthodox religion.

To sum up, I'll just give two quotes; the first on the extreme difficulty of following the way to God, and the second on the different paths that lead to the same end:

Enter ye in at the strait gate: for wide is the gate, and broad is the way, that leadeth to destruction, and many there be which go in thereat: Because strait is the gate, and narrow is the way, which leadeth unto life, and few there be that find it. -- Matthew 17:13-14

With regard to the highest human goal, there is no contradiction among scriptures, since all teach the very same reward: deliverance. Nevertheless, differing salvific paths are taught, according to the intellect of the beings to be favored. This omniscient Lord taught various kinds of approaches when he saw that ``these people can be helped to reach beatitude in the way they prefer on this path.'' Just as people from a crowd that wants to enter a single fort or a big house also enter through different doors, liberation-seekers too enter the highest abode in the same way. The following wise saying of Jayanta, the prodigy, who has mastered the essence of all sciences, who knows reality, and who has shaken off error, refers to the same thing:

``The many means taught by various scriptural approaches converge in the single summum bonum, as the currents of the Ganges meet in the ocean.'' -- A?gamad?ambara 4110-4115

\hfill

\texttt{me on 2016-05-05 at 14:24 said:}

Thank you very much for your comprehensive reply. I will be mulling it over.

\end{sffamily}
\end{footnotesize}

